% Fragment: the seed leg of pass^k under an opt-in execution-noise model.
% \input from 05-experiments.tex after the relative-mandate fragment, or cite
% from the Reliability paragraph of 03-benchmark.tex.
% Evidence: paper/evidence/final/seed-leg.jsonl (2646 records: 2592 per-run
% rows and 54 per-agent verdict rows, produced by the command recorded on
% every record: cargo run --release -p sharpebench-harness --example
% seed_leg_eval -- paper/evidence/final/seed-leg.jsonl).

\subsection{The seed leg under a richer execution-noise model}\label{sec:seedleg}

The Reliability paragraph of \cref{sec:benchmark} concedes that the seed leg of $\text{pass}^{k}$ is a narrow stability check: under the shipped cost model the eight execution seeds vary only the base slippage draw, a few basis points, so every refusal on the paper's grid is attributable to the window leg. This section adds an opt-in execution-noise model under which the seeds resample fills rather than basis points, re-runs the reference field under both profiles on three daily datasets, and answers two questions the concession leaves open: whether the seed leg ever becomes the refusing leg once execution is genuinely noisy, and how far per-seed Sharpes then disperse.

\paragraph{Definition.} Let an order at bar $t$ on symbol $i$ under execution seed $s$ request a change $\Delta$ in position value toward its target weight, after the liquidity cap. The model draws three uniforms $u_d, u_f, u_q \sim \mathcal{U}[0,1)$ from a SplitMix64 stream keyed on the triple, $R(s, t, i)$, so each quantity below is a pure function of $(s, t, i)$ and the base slippage draw of the default model is untouched. Three effects apply in order.
\begin{align}
\text{fill bar}\quad \tau &= t + \mathbb{1}[\,u_d < p_d\,],\label{eq:noise-delay}\\
\text{filled value}\quad \Delta_{\rm fill} &= \begin{cases}\varphi\,\Delta & \text{if } (1-\varphi)\,|\Delta| > \kappa\,\mathrm{NAV}_t,\\ \Delta & \text{otherwise,}\end{cases}\qquad \varphi = \varphi_{\min} + (1-\varphi_{\min})\,u_f,\label{eq:noise-partial}\\
\text{queue slippage}\quad s_q &= u_q\;\rho_t\;\min\!\Big(1,\ \frac{\pi}{\pi_{\rm ref}}\Big),\qquad \rho_t = \Big|\frac{c_t}{c_{t-1}} - 1\Big|,\quad \pi = \frac{|\Delta_{\rm fill}|}{\mathrm{NAV}_t}.\label{eq:noise-queue}
\end{align}
A fresh order uses \cref{eq:noise-delay}; a carried order fills at its scheduled bar without a second delay. A delayed order and the unfilled remainder $\Delta - \Delta_{\rm fill}$ of a partial fill (\cref{eq:noise-partial}) are carried to bar $t+1$ as an order for the same target weight, where they fill at bar $t+1$'s price before that bar's fresh orders; a fresh order on the same symbol supersedes the carried one, which is cancel-and-replace. The carry floor $\kappa$ makes a remainder worth less than $\kappa$ of NAV fill in full, so a carry drains in finitely many bars instead of shrinking geometrically forever. The queue term (\cref{eq:noise-queue}) is an adverse move of up to one bar's range, scaled linearly by participation up to $\pi_{\rm ref}$, and is added to the base seeded slippage and the square-root impact term before the fill price is formed, with the sign against the trade as before. The frozen datasets carry closes only, so $\rho_t$ is the close-to-close absolute move rather than a high-low range; it is a proxy, and \cref{eq:noise-queue} should be read as such. The shipped parameters are $p_d = 0.25$, $\varphi_{\min} = 0.5$, $\kappa = 10^{-3}$ and $\pi_{\rm ref} = 0.10$; none is calibrated to a venue, and the model is stylized by design.

The model is opt-in. \texttt{CostModel} gains a \texttt{noise: Option<ExecutionNoise>} field that deserializes to \texttt{None} from every existing configuration, and with \texttt{None} the fill arithmetic is the historical path in the same order, so the golden fixtures, the native/wasm parity tests and the evidence sweeps are byte-identical. \texttt{CostProfile::Realistic} resolves to the typical cost model with \texttt{ExecutionNoise::default()}. The carried orders live in the book state and are absent from a serialized snapshot when empty. Given $(s, t, i)$ every quantity is deterministic, and two runs under the same seed replay bit for bit.

\paragraph{Evidence.} The field of \cref{sec:riskmanaged} (buy-and-hold, momentum, hold, the risk-managed agent, and five luck-floor agents), the same window rule and the same eight execution seeds, on three daily datasets from three market panels: US indices (2513 bars, 3 symbols, six windows of 408 bars, regimes UUUDUU), crypto majors (1000 bars, 5 symbols, six windows of 156, regimes UCUUDD) and FX majors (4157 bars, 5 symbols, six windows of 682, regimes UDCUDU), scored under the default \texttt{ScoreConfig} for each dataset's periods per year. Every (window, seed) cell is tested against the per-run bar $\text{PSR}(r, 0) \ge 0.90$ through \texttt{sharpebench\_core::per\_run\_passes}, and each agent's verdict is decomposed by leg: a window whose eight sampled execution draws all fail is a \emph{window-leg} refusal, while a window with some sampled draws passing and some failing is a \emph{seed-leg} refusal. The decomposition is exhaustive for this eight-draw experiment: $\text{pass}^{k}$ holds if and only if neither leg refuses. For each (agent, window) the across-seed sample standard deviation of per-run annualized Sharpe is recorded, and for scale the across-window standard deviation of the per-window mean Sharpe. \Cref{tab:seedleg} lists the three trading reference agents on every dataset under both profiles; hold never trades, so both profiles give it identically zero returns, and the five luck-floor agents are summarized in the text because their seed dispersion is not an execution effect.

\begin{table}[t]
\centering
\footnotesize
\setlength{\tabcolsep}{3.5pt}
\caption{The seed leg of $\text{pass}^{k}$ under the default cost model and under \texttt{CostProfile::Realistic} (\cref{eq:noise-delay,eq:noise-partial,eq:noise-queue}), three daily datasets, eight seeds, six windows. Windows: P when all eight seeds pass the per-run bar, M when the seeds disagree, a period when all eight fail. Leg is the refusing leg. Seed std is the across-seed sample standard deviation of per-run annualized Sharpe, mean and maximum over the six windows; window std is the across-window standard deviation of the per-window mean Sharpe; mean SR is the mean per-run annualized Sharpe over the 48 cells. Closest fail is the highest PSR of any cell in an all-fail window, the nearest the window leg came to admitting a window. No agent passes $\text{pass}^{k}$ or is rank-eligible under either profile.}
\label{tab:seedleg}
\begin{tabular}{lllcccccrr}
\toprule
Dataset & Agent & Profile & Windows & Leg & \multicolumn{2}{c}{Seed std} & Window std & Mean SR & Closest fail \\
 & & & (0 to 5) & & mean & max & & & \\
\midrule
US indices 1d & buy-and-hold & default   & \texttt{P...P.} & windows & 0.0001 & 0.0002 & 0.8153 & $+0.984$ & 0.879 \\
US indices 1d & buy-and-hold & realistic & \texttt{P...P.} & windows & 0.0063 & 0.0092 & 0.8073 & $+0.966$ & 0.878 \\
US indices 1d & momentum     & default   & \texttt{......} & windows & 0.0025 & 0.0034 & 0.8849 & $-0.535$ & 0.723 \\
US indices 1d & momentum     & realistic & \texttt{......} & windows & 0.2210 & 0.3142 & 0.9366 & $-0.897$ & 0.674 \\
US indices 1d & risk-managed & default   & \texttt{......} & windows & 0.0019 & 0.0033 & 0.8637 & $-0.163$ & 0.861 \\
US indices 1d & risk-managed & realistic & \texttt{......} & windows & 0.2079 & 0.3446 & 0.8081 & $-0.801$ & 0.688 \\
\midrule
Crypto 1d     & buy-and-hold & default   & \texttt{P..P..} & windows & 0.0001 & 0.0001 & 1.7941 & $+0.716$ & 0.890 \\
Crypto 1d     & buy-and-hold & realistic & \texttt{P..P..} & windows & 0.0332 & 0.0536 & 1.8182 & $+0.679$ & 0.882 \\
Crypto 1d     & momentum     & default   & \texttt{......} & windows & 0.0019 & 0.0024 & 1.7856 & $-0.631$ & 0.898 \\
Crypto 1d     & momentum     & realistic & \texttt{......} & windows & 0.4354 & 0.5969 & 2.0989 & $-1.982$ & 0.687 \\
Crypto 1d     & risk-managed & default   & \texttt{P..P..} & windows & 0.0009 & 0.0013 & 1.6159 & $+0.574$ & 0.854 \\
Crypto 1d     & risk-managed & realistic & \texttt{P.....} & windows & 0.2579 & 0.4267 & 1.9319 & $-0.311$ & 0.785 \\
\midrule
FX majors 1d  & buy-and-hold & default   & \texttt{......} & windows & 0.0001 & 0.0001 & 0.6502 & $-0.070$ & 0.774 \\
FX majors 1d  & buy-and-hold & realistic & \texttt{......} & windows & 0.0074 & 0.0114 & 0.6442 & $-0.084$ & 0.775 \\
FX majors 1d  & momentum     & default   & \texttt{......} & windows & 0.0034 & 0.0044 & 0.8563 & $-4.142$ & 0.000 \\
FX majors 1d  & momentum     & realistic & \texttt{......} & windows & 0.1742 & 0.2265 & 0.7068 & $-4.031$ & 0.000 \\
FX majors 1d  & risk-managed & default   & \texttt{......} & windows & 0.0036 & 0.0109 & 0.5786 & $-1.917$ & 0.035 \\
FX majors 1d  & risk-managed & realistic & \texttt{......} & windows & 0.1263 & 0.1745 & 0.5882 & $-2.345$ & 0.033 \\
\bottomrule
\end{tabular}
\end{table}

\paragraph{Result.} The realistic profile makes the seed leg measure something, and it still never refuses. Pooled over the three trading agents and six windows, the across-seed standard deviation of per-run annualized Sharpe rises from $0.0015$, $0.0010$ and $0.0024$ under the default profile on US indices, crypto and FX to $0.1451$, $0.2422$ and $0.1026$ under the realistic profile, roughly two orders of magnitude; the largest single-window value is $0.5969$ (momentum on crypto). For the two agents that trade every bar the ratio of per-window seed dispersion to across-window dispersion is now between $0.13$ and $0.26$, where under the default it was below one percent; for buy-and-hold, which trades once and then drifts, the seed dispersion stays below $0.06$ because the model charges the fills, and buy-and-hold has almost none. The refusing leg is nevertheless \emph{windows} in all eighteen rows of \cref{tab:seedleg} under both profiles: for the reference agents not one window is mixed. The one change in a window verdict, risk-managed on crypto losing its second passing window (\texttt{P..P..} to \texttt{P.....}), is a whole-window shift, all eight seeds failing where all eight had passed, driven by the level effect of the noise on a rebalancing agent (mean Sharpe from $+0.574$ to $-0.311$), not by a seed-dependent verdict. The level effect is itself the largest visible consequence: momentum loses $0.36$, $1.35$ and gains $0.11$ of annualized Sharpe on the three datasets, risk-managed loses $0.64$, $0.88$ and $0.43$, and buy-and-hold loses $0.02$, $0.04$ and $0.01$. Among the luck-floor agents the two mixed windows that exist on the grid (luck-floor-00 and luck-floor-01 on crypto under the default profile, two and one of eight seeds passing window 0) disappear under the realistic profile, where every luck-floor window fails on all eight seeds; but the luck floor's seed dispersion of $0.14$ to $0.31$ under the default profile is a policy effect, each \texttt{RandomAgent} being re-seeded per run, and says nothing about execution noise.

\paragraph{Interpretation.} The seed leg does not refuse because the per-window PSRs of these agents sit far from the bar on both sides. The Closest fail column shows the nearest any all-fail window came to $0.90$: $0.898$ for momentum on crypto under the default, where the across-seed PSR range on that window was $0.0014$, and $0.687$ under the realistic profile, where the range had grown to $0.464$ but the level had moved away from the bar faster than the dispersion grew. On the passing side the tightest margin is buy-and-hold's second passing window on crypto, whose lowest seed PSR is $0.9012$ under the default and $0.9086$ under the realistic profile, with the across-seed PSR range on that window at $0.0001$ and $0.0402$; the realistic profile brings the seed range to the same order as that margin, and the window stays unanimous by roughly one range width. That is the honest reading: under the shipped model the seed leg cannot flip a window because the seeds barely move the PSR, and under the realistic model it could flip a window that sits within a few hundredths of the bar, but on this grid no window does. The paper's statement that every refusal is a window-leg refusal therefore survives the richer noise model on three panels, with the qualification that the seed leg is now a live test rather than a formality.

Three caveats bound the claim. The noise model is stylized: its four parameters are set by hand, the range proxy is close-to-close because the datasets carry no high or low, and carried orders fill at the next close, so the model has the structure of delay, partial fill and queue loss but not the magnitudes of any particular venue. The across-seed standard deviations are estimated from eight seeds, for which the relative standard error of a sample standard deviation is about $27$ percent under an IID normal approximation (\cref{sec:incentives}), so the two-orders-of-magnitude contrast is robust and the third digit is not. And three daily datasets from three panels are enough to show that the seed leg can be made to measure execution and still not be the refusing leg for these agents, not enough to say it never is; an agent whose per-window Sharpe sits near the bar, which none of the author-written references does, is exactly the case in which the realistic seed leg would decide, and that is the case the arena is meant to see.
